%%=============================================================================
%% Inleiding
%%=============================================================================

\chapter{\IfLanguageName{dutch}{Inleiding}{Introduction}}%
\label{ch:inleiding}


Modern society relies heavily, we could even say entirely, on the permanent availabiltiy of our communication netwroks.
Everything has switched to global networks, the Internet, cellular data, phone networks, etc. Everything became embedded in day-to-day life that people use it without even thinking about it, we take it for granted.
Everything means even emergency response depends on said networks to communicate and to coordinate first responders, warn people and keep citizens connected.
This means that this dependcy can becoime critical at the exact moment it is most needed, when networks become unavailable.

When a big storm destroys a cellular tower, or when floods cut power of a specific region (where bases might be running) or other large-scale incidents happen that infrastructure just disappears, and suddenly what we took for granted is no longer available.
This is not just hypotetical; natural disasters, power outages, conflict situations (disrupting traffic, censoring, etc.), these are moments when using the network is necessary, but impossible.

For those exact situations a back-up system should be available. Not only for first responders or civil protection but also for the general public.
A fallback available to the public should be quick and easy to deploy, not require a specialist and be affordable enough.
Several technologies already exists, but they all have their limiations which make them difficult to effectively applu in an emergency.
These limitations can be cost, licences, needed hardware, extra fees, etc.

Thios bachelor thesis will investigate \textbf{Meshtastic}. \textbf{Meshtastic} is an open-source project that allows to build a decentralized, encrypted mesh network. This can be done on low-cost \gls{LoRa} radio hardware.
Meshtastic doesn't require any SIM card, internet connection or lcience. It provides self-healing multi-hop routing which allows us to extend the initial coverage without needing any fixed existing infrastructure.
For these exact reasons, Meshtastic might be a promising candidate for the emergency contexts described above. However real-world tests have been conducted a real emergency reliability and emergency hasn't been evaluated yet. This thesis will address that gap in research by creating a PoC in the region around Ronse, Belgium.
\section{\IfLanguageName{dutch}{Probleemstelling}{Problem Statement}}%

\label{sec:probleemstelling}

Whenever the main infrastructure fails, the responsible people for coordinating instructions, help, etc. are left without any way to do this.
The core problem of this thesis if the absence of an accessible, affordable, reliable and easy to deploy solution that functions independtly of the failed infraastrructure.
Also no evidence is available about how well Meshtastic handles this role in practice.

The primary audience are IT professionals that are responsible for continuity and civil protectiuon communication, such as security teams etc.
As a secondary audience we could also identify volunteer oprganizations, municvipal safety services and network engiuveers who want to investigate low-sot, off the grid coimmunication options.
A last additional audience could just be general public, for who privacy matters and don't want to roam on a controlled network.

\section{\IfLanguageName{dutch}{Onderzoeksvraag}{Research question}}%
\label{sec:onderzoeksvraag}

The central research question of the thesis is: 
\begin{quote}
\emph{To what extent does Meshtastic, as a decentralized \gls{LoRa}-based mesh network, offer a reliable and readily
deployable emergency communication solution when primary network infrastructure fails?}
\end{quote}

From this main question two domains can be identified: the problem and solution domain.

The \textbf{problem domain} is examining the requirements that occur whenever infrastructure fails and how existing solutions address these.
To do this following sub-questions will be answered:
\begin{itemize}
  \item What gaps exist in emergency contexts when the primary infra becomes unavailable?
  \item How do existing alternatives compare to Meshtastic in terms of accessibility, cost, required expertise, etc.?
  \item What environmental factors of the Belgian (Flemish Ardennes and Wallonia) landscape and terrain affect LoRa performance once deployed?
  \item How secure is the Meshtastic encryption in case of traffic interception?
\end{itemize}

The \textbf{solution domain} on the other hand is evaluating the measurements from the real-life deplooyement during the PoC.
This provides us these sub-questions:
\begin{itemize}
  \item What range and reliability does a Meshtastic network have across different topologies (direct long-distance
  vs.\ multi-hop) in a real deployment (East Flanders/Hainaut province: Flemish Ardennes/Pays des Collines regions (Ronse/Oudenaarde/Kluisbergen/Tournai))?
  \item How does the Meshtastic security model hold up against potential attacks?
  \item How does it compare with other alternatives in terms of requirements in emergency situations?
\end{itemize}

\section{\IfLanguageName{dutch}{Onderzoeksdoelstelling}{Research objective}}%
\label{sec:onderzoeksdoelstelling}

There are two goals for this thesis.
First off, using literature, comparing Meshtastic to the most relevant alternatives (being APRS, DMR and other commercial satelitte messengers).
The second one, making use of the \gls{poc}, measuring Meshtastic's effective range, reliability and deployabvility in a real-life mixed terrain.

What the concrete deliverables are:
\begin{itemize}
  \item A requirements list and comparison matrix
  \item A working three or four node Meshtastic deployment (fixed base node, fixed relay node(s) and a mobile node)
  \item A report of field measurements collected during PoC test sessions (range, \gls{rssi}, \gls{snr}, \gls{pdr} and hop count)
  \item A final matrix comparing Meshtastic to the alternatives
\end{itemize}

Success for this thesis can be defined as being able to give a clear, evidence-based answer to the research question, sup^ported by th efield measuremnts the PoC provided.

\section{\IfLanguageName{dutch}{Opzet van deze bachelorproef}{Structure of this bachelor thesis}}%
\label{sec:opzet-bachelorproef}

The structure of this thesis is as follows:

Chapter~\ref{ch:inleiding} introduces the main problem, questions and ideas.

%TODO